\thispagestyle{plain}
\begin{center}
{\zihao{2}\sffamily\bfseries\color{DesertNavy} 基于时空网络、随机动态规划与拥挤博弈的穿越沙漠策略研究\par}
\vspace{0.6em}
{\color{DesertNavy}\rule{0.38\textwidth}{1.2pt}}\par
\vspace{1.2em}
{\zihao{3}\sffamily\bfseries 摘\quad 要}
\end{center}

\begingroup
\small
本文研究有限资金、承重与天气约束下的沙漠穿越问题。六关的核心差异并非地图大小，而是信息集由“未来天气全知”依次变化为“仅知当天气象”和“可观察其他玩家历史”。据此，本文建立“时间扩展网络混合整数规划—有限时域马尔可夫决策过程—多人动态博弈”的递进框架，并用独立规则模拟器、全序列枚举、预先规定风险预算和最佳响应检验形成分级证据链。

针对问题一，将位置流、逐日行动、资源库存、现金、负重、村庄补给与矿山收益统一写入时间扩展网络 MILP。HiGHS 对第一、二关均返回全局最优，独立模拟器逐日回放后，第一关以水178箱、食物333箱出发，第24天到达，最终资金为\textbf{10470元}；第二关以水130箱、食物405箱出发，第30天到达，最终资金为\textbf{12730元}。回放还发现公开参考表第二关第15--18天水量少扣10箱，但修正后最优值不变，说明“数值最优”和“逐日合法”必须分别审计。

针对问题二，题面未给未来天气分布，故最佳策略只能是条件结论。第三关在晴朗/高温独立分布下建立天气先揭示的精确有限时域 MDP；高温概率为0.2、0.5、0.8时，最优初始装载均为水54箱、食物54箱，沿 $1\to4\to6\to13$ 三日直达，期望终值分别为9382、9310、9238元，并对每档全部 $2^{10}=1024$ 条天气序列实现零失败。第四关先在IID下从567个装载—沙暴预算候选选出240/240箱、B5效率策略；再构造$P_\rho=\rho I+(1-\rho)\bm 1p^\top$，在保持天气边际概率不变的条件下把持续性扩展到$\rho\le0.35$。分布鲁棒重选得到240/240箱、B9策略；在最不利边际、$\rho=0.35$的30000条样本中，失败率为0.1967\%（Wilson上界0.2536\%），而IID-B5为2.2300\%。

针对问题三，首先从原始 Word 公式核实同边 $k$ 人移动的消耗系数为 $2k$、同矿收益为基础收益的 $1/k$。第五关对固定对手计划执行完整最佳响应动态规划，得到终值\textbf{9535元}和\textbf{9510元}的非对称纯策略 Nash 均衡，双方最大单边偏离增益均为0。第六关实现三人联合反馈模拟；机械复制单人挖矿策略在IID基准天气下玩家失败率高达97.865\%。进一步按角色校准二维装载，最终采用215/215、215/215与225/230箱的异路错峰剖面；在最不利边际、$\rho=0.35$下玩家失败率0.3400\%、固定角色最大Wilson上界0.4344\%，总体均值\textbf{7383.3元}。受约束最佳响应预言机对每名玩家完整检查70条最短路乘2种风险可行错峰装载，共140个偏离，样本外最大可利用度为\textbf{0元}。公开随机分配三种角色实现事前公平。

本文的创新点在于：以信息结构而非题号选择模型；将失败样本按0元纳入总体效用以消除幸存者偏差；用保持边际不变的Markov模型分离天气频率与时间聚集；将“全局最优、精确枚举、均衡偏离与分布鲁棒样本外推荐”严格分级。$\rho=0.70$的集合外破坏性测试明确展示保护边界，因而不把条件安全包装成任意天气下的无条件最优。

\vspace{0.8em}
\noindent\textbf{关键词：}时间扩展网络；混合整数规划；马尔可夫决策过程；鲁棒策略；Nash均衡；拥挤外部性；样本外验证
\endgroup
